\documentclass[11pt,letterpaper]{article}

\usepackage[margin=1in]{geometry}
\usepackage[T1]{fontenc}
\usepackage[utf8]{inputenc}
\usepackage{lmodern}
\usepackage{microtype}
\usepackage{amsmath,amssymb,bm}
\usepackage{siunitx}
\usepackage{booktabs}
\usepackage{tabularx}
\usepackage{longtable}
\usepackage{enumitem}
\usepackage{xcolor}
\usepackage{hyperref}

\hypersetup{
    colorlinks=true,
    linkcolor=black,
    citecolor=blue!60!black,
    urlcolor=blue!60!black
}

\title{\textbf{Phase 0 Computational Setup and Modeling Strategy}}
\author{Daniel Alejandro Olaya}
\date{\today}

\begin{document}
\maketitle

\section{Phase 0 Objective}

Phase 0 is designed to compare three representative baseline architectures under controlled and identical simulation conditions: an organic/biomimetic tubular structure, a hybrid organic--inorganic tubular structure, and a predominantly inorganic tubular structure. The objective is to identify which baseline architecture maximizes physical plausibility, computational tractability, and measurable output contrast.

\section{Computational Environment}

The initial computational setup is developed on macOS. This environment will be used for documentation, geometry generation, small test systems, scripting, post-processing, plotting, and visualization. Larger production simulations may later be migrated to a Linux workstation, HPC system, or cloud compute environment.

\section{Software Stack and Justification}

The initial computational environment is a macOS workstation used for documentation, model generation, local testing, trajectory analysis, and plotting. Larger production simulations may be migrated later to a Linux workstation, HPC system, or cloud compute environment depending on system size, number of replicas, trajectory length, and computational cost.

\subsection{LAMMPS}

LAMMPS is available locally through the executable \texttt{lmp\_mpi}. The installed version is LAMMPS 10 Dec 2025, compiled with Open MPI v5.0.7. The build includes several packages relevant to this project, including \texttt{MOLECULE}, \texttt{KSPACE}, \texttt{DIPOLE}, \texttt{DRUDE}, \texttt{DIELECTRIC}, \texttt{MANYBODY}, \texttt{REAXFF}, \texttt{QEQ}, \texttt{RIGID}, \texttt{SPIN}, \texttt{PHONON}, \texttt{BROWNIAN}, \texttt{COLLOID}, \texttt{CORESHELL}, and \texttt{DPD}.

LAMMPS is selected as the primary engine for inorganic, hybrid organic--inorganic, coarse-grained, and field-driven models because it provides broad support for particle-based materials modeling, long-range electrostatics, dipolar interactions, external fields, rigid bodies, many-body potentials, reactive force fields, and specialized packages for dielectric, spin, and phonon-related models. Its flexibility is particularly relevant for Phase 0 because the candidate architectures include inorganic tubes, hybrid coated tubes, and confined polar media. The use of LAMMPS for large-scale particle-based molecular simulation is supported by the foundational and modern LAMMPS literature.

The current macOS LAMMPS build is suitable for local setup, geometry validation, debugging, and small-to-medium test systems. Because OpenMP is not enabled at the compiler level and no GPU acceleration is reported, large production simulations may be more efficiently executed later on a Linux workstation or HPC/cloud environment.

\subsection{GROMACS}

GROMACS is available locally through the executable \texttt{gmx}. The installed version is GROMACS 2025.4, compiled in mixed precision with thread-MPI, OpenMP support, AVX2 SIMD instructions, and FFTW. GPU support is disabled in the local build.

GROMACS is selected for biomolecular, peptide-like, organic, and soft-matter candidates because it is highly optimized for classical molecular dynamics of water-containing molecular systems and provides mature tools for energy minimization, equilibration, trajectory generation, and standard structural analyses. It is particularly suitable if the organic/biomimetic Phase 0 candidate is represented as a peptide-like or biomolecular tubular assembly. The use of GROMACS for efficient molecular simulations is supported by standard GROMACS references.

The current macOS GROMACS build is suitable for local testing and moderate CPU-based simulations. Large-scale production runs may benefit from optimized Linux/GPU builds.

\subsection{Python Analysis Environment}

Python 3.12.12 is available through the Miniforge environment. Python will be used for geometry generation, workflow automation, trajectory analysis, dipole-orientation metrics, autocorrelation functions, plotting, tabular summaries, and preparation of report-ready figures. The initial package list includes \texttt{NumPy}, \texttt{SciPy}, \texttt{pandas}, \texttt{matplotlib}, \texttt{MDAnalysis}, \texttt{MDTraj}, \texttt{Jupyter}, and \texttt{tqdm}.

\subsection{Later-Stage Tools}

Quantum ESPRESSO and Octopus are reserved for later electronic-structure and optical-response modules, including DFT, TDDFT, transition-dipole calculations, polarizability estimates, and fragment-level electronic parameterization. COMSOL is reserved for later continuum electromagnetic and device-level modeling once effective dielectric, polarizability, or conductivity parameters are available from molecular simulations, electronic calculations, or literature data.

\section{Repository Structure}

The repository is organized to separate documentation, references, scripts, systems, controls, simulations, results, figures, logs, and daily notes.

\section{Candidate Architecture Definition}

The first validated geometric reference is a tubular scaffold with outer diameter
\SI{24.0}{nm}, lumen diameter \SI{14.0}{nm}, wall thickness \SI{5.0}{nm}, and
segment length \SI{20.0}{nm}. This geometry is used as a common reference for
Phase 0 so that the organic/biomimetic, hybrid organic--inorganic, and
predominantly inorganic candidates can be compared under identical dimensional
constraints.

The initial generated scaffold is a geometric point-cloud representation rather
than a force-field-ready molecular model. Its purpose is to verify dimensions,
coordinate conventions, plotting routines, and documentation workflow before
constructing chemically specific candidate systems. Geometry inspection confirmed
an estimated lumen diameter of \SI{14.0}{nm}, outer diameter of \SI{24.0}{nm},
segment length of \SI{20.0}{nm}, and 17,280 shell points.

\begin{table}[h!]
\centering
\caption{Initial Phase 0 candidate matrix.}
\begin{tabularx}{\textwidth}{p{3.1cm} X X}
\toprule
\textbf{Class} & \textbf{Initial candidate} & \textbf{Primary purpose} \\
\midrule
Organic / biomimetic &
Peptide-like or aromatic peptide tubular scaffold with polar wall groups &
Test biomimetic confinement, hydration-layer behavior, polar wall effects, and structural persistence. \\
Hybrid organic--inorganic &
BNNT or CNT scaffold with polar or peptide-like coating &
Test whether a mechanically robust inorganic tube plus polar coating improves confined-water ordering and dielectric response. \\
Predominantly inorganic &
BNNT with confined water &
Test robust tubular confinement, confined-water ordering, dipole autocorrelation, and field-induced dielectric response. \\
\bottomrule
\end{tabularx}
\end{table}

\section{Control Systems}

The Phase 0 comparison requires controls to determine whether the measured
response is caused by structured confinement, confined water, wall chemistry, or
ordered functionalization. The initial controls are:

\begin{itemize}
    \item dry or empty tube, to isolate the role of confined water;
    \item water-filled unfunctionalized tube, to separate confinement from wall functionalization;
    \item bulk or unconfined water, to compare confined and unconfined relaxation behavior;
    \item randomized wall polarity, when applicable, to test whether ordered dipolar patterning matters.
\end{itemize}

\section{Simulation Protocol}

Pending definition.

\section{Analysis Metrics}

Pending definition.

\section{Results}

Pending simulations.

\section{Discussion}

Pending results.

\section{Next-Step Decision Criteria}

The first decision criterion is whether the computational setup is sufficiently reproducible and flexible to support the Phase 0 comparison.

\end{document}
